%!TEX root = ../report.tex
\documentclass[../report.tex]{subfiles}

\begin{document}
    \section{Executing the Code}
    \label{sec:executing_code}

    % Start with describing the environment and then describe how you executed the code in the simulation and on the real robot.
    % this section is to show the results of the executed code.

    We first tested all tasks in the simulation. On the real robot, we tested the exploration (task 3).
    This section describes both environments and shows the results.

    \subsection{The Simulation}
    \label{sec:executing_code:simulation}

    % dont forget to add pictures and videos of the simulation

    \textbf{Environment.}
    We used the Robile simulation in Gazebo Classic with ROS~2 Humble.
    It ran in Docker containers: partly on an Ubuntu computer in a container with Ubuntu, and mostly on a Mac, where the container emulates the amd64 architecture, because Gazebo Classic is not available for ARM processors.
    As the Gazebo interface did not open in the container, we only ran the simulation server and used RViz to see the robot, the laser scan, the map and the particles.
    Most tests were done in the world \emph{closed\_walls} (about \qtyproduct{7 x 7}{m}) and some in a house world as a harder test with open doors.

    \textbf{Task 1: Path and motion planning.}
    For this task, the robot was localised with our particle filter in a map created before.
    We sent three goals one after the other (Fig.~\ref{fig:sim_task1}).
    The robot reached all goals without collisions in about \SIrange{30}{50}{s} each, with a final position error of about \SIrange{0.02}{0.16}{m}.
    With the new potential field in the configuration space (Sec.~\ref{sec:challenges:real_robot}), the robot also reached all three goals, with an error of about \SIrange{0.03}{0.14}{m}.
    Video: \href{https://youtu.be/eph5LCBV-MY}{aufgabe1 navigation}

    \begin{figure*}[t]
        \centering
        \begin{subfigure}{0.32\linewidth}
            \includegraphics[width=\linewidth]{figures/sim/task1_a.png}
            \caption{Start}
        \end{subfigure}
        \hfill
        \begin{subfigure}{0.32\linewidth}
            \includegraphics[width=\linewidth]{figures/sim/task1_b.png}
            \caption{On the way to a goal}
        \end{subfigure}
        \hfill
        \begin{subfigure}{0.32\linewidth}
            \includegraphics[width=\linewidth]{figures/sim/task1_c.png}
            \caption{Goal reached}
        \end{subfigure}
        \caption{Task 1 in the simulation (RViz): The robot follows the waypoints of the A* path to the goal (purple arrow).}
        \label{fig:sim_task1}
    \end{figure*}

    \textbf{Task 2: Localisation.}
    We started the particle filter without a start pose, so the particles were spread over the whole map (Fig.~\ref{fig:sim_task2}).
    While the robot was moving, the particles converged to the correct pose after about \SIrange{15}{90}{s}, depending on the start position.
    After that, the error of the estimated position was about \SIrange{0.01}{0.15}{m}.
    In some runs, the filter first converged to a wrong but similar-looking place and corrected itself after \SIrange{30}{80}{s} thanks to the random particles.
    Video: \href{https://youtu.be/jE-axqhAQd0}{aufgabe2 lokalisierung}

    \begin{figure*}[t]
        \centering
        \begin{subfigure}{0.32\linewidth}
            \includegraphics[width=\linewidth]{figures/sim/task2_a.png}
            \caption{Particles spread over the map}
        \end{subfigure}
        \hfill
        \begin{subfigure}{0.32\linewidth}
            \includegraphics[width=\linewidth]{figures/sim/task2_b.png}
            \caption{Robot is moving}
        \end{subfigure}
        \hfill
        \begin{subfigure}{0.32\linewidth}
            \includegraphics[width=\linewidth]{figures/sim/task2_c.png}
            \caption{Particles converged}
        \end{subfigure}
        \caption{Task 2 in the simulation: global localisation with the particle filter (red arrows are particles).}
        \label{fig:sim_task2}
    \end{figure*}

    \textbf{Task 3: Environment exploration.}
    The robot started in an unknown environment and explored it on its own (Fig.~\ref{fig:sim_task3}).
    In the final runs in \emph{closed\_walls}, the complete map (about \SI{52}{m^2}) was created in about \SIrange{90}{120}{s}, and the exploration stopped by itself when no reachable fringe was left.
    Video: \href{https://youtu.be/DoE1StjPe04}{aufgabe3 exploration}

    \begin{figure*}[t]
        \centering
        \begin{subfigure}{0.32\linewidth}
            \includegraphics[width=\linewidth]{figures/sim/task3_a.png}
            \caption{Start}
        \end{subfigure}
        \hfill
        \begin{subfigure}{0.32\linewidth}
            \includegraphics[width=\linewidth]{figures/sim/task3_b.png}
            \caption{During the exploration}
        \end{subfigure}
        \hfill
        \begin{subfigure}{0.32\linewidth}
            \includegraphics[width=\linewidth]{figures/sim/task3_c.png}
            \caption{Final map}
        \end{subfigure}
        \caption{Task 3 in the simulation: The map created by gmapping grows while the explorer selects goals at the map fringe.}
        \label{fig:sim_task3}
    \end{figure*}

    \subsection{The Real Robot}
    \label{sec:executing_code:real_robot}

    \textbf{Environment.}
    In the lab, we used a Robile with four wheels and a laser scanner at its front.
    The ROS~2 nodes ran in a Docker container with Ubuntu on a laptop with Ubuntu, which was connected to the robot via WLAN.
    We built an arena in our lab room (Fig.~\ref{fig:real_arena}): an area of about \qtyproduct{4 x 6}{m}, bounded by boxes, cupboards, shelves and a sofa, with a narrow passage between two parts of the area.
    Compared to the simulation, we reduced the speed of the robot to at most \SI{0.3}{m/s} and \SI{0.8}{rad/s}.

    \begin{figure}[ht]
        \centering
        \includegraphics[width=0.9\linewidth]{figures/arena.jpg}
        \caption{The Robile in our arena in the lab room, bounded by boxes and furniture.}
        \label{fig:real_arena}
    \end{figure}

    \textbf{Real world explorer.}
    The robot started in the arena without a map and explored it with the explorer and \emph{slam\_gmapping} (Fig.~\ref{fig:real_task3}).
    In the first tests, the robot did not pass the narrow passage to the second part of the arena (see Sec.~\ref{sec:challenges:real_robot}).
    With the new navigator and the improved explorer, the robot drove through the narrow passage and mapped the whole arena.
    Videos: \href{https://youtu.be/7wGQjqQ4EhQ}{Mapping in Lab; Real Robot}, \href{https://youtu.be/zNW6HkEr-Kk}{Mapping in Lab; Screenrecording}

    \begin{figure}[ht]
        \centering
        \includegraphics[width=0.6\linewidth, trim=150 160 120 65, clip]{figures/lab_map-5.png}
        \caption{Map of the arena created by the real world explorer (white: free, black: occupied, grey: unknown).}
        \label{fig:real_task3}
    \end{figure}

\end{document}
